% ══════════════════════════════════════════════════════════════════════════════
% KTD.jl Technical Report v0.2
% Rod Read — Windswept & Interesting Ltd
% July 2026
% ══════════════════════════════════════════════════════════════════════════════

\documentclass[11pt,a4paper]{article}

% ── Typography ───────────────────────────────────────────────────────────────
\usepackage[T1]{fontenc}
\usepackage[utf8]{inputenc}
\usepackage{mathptmx}
\usepackage{microtype}
\usepackage{setspace}
\onehalfspacing

% ── Page geometry ────────────────────────────────────────────────────────────
\usepackage[top=3cm,bottom=3cm,left=3.2cm,right=3.2cm]{geometry}

% ── Headers / footers ────────────────────────────────────────────────────────
\usepackage{fancyhdr}
\pagestyle{fancy}
\fancyhf{}
\fancyhead[L]{\small\textit{KTD.jl Technical Report v0.2 --- July 2026}}
\fancyhead[R]{\small\textit{Windswept \& Interesting Ltd}}
\fancyfoot[C]{\thepage}
\renewcommand{\headrulewidth}{0.4pt}

% ── Figures & graphics ───────────────────────────────────────────────────────
\usepackage{graphicx}
\usepackage[font=small,labelfont=bf,labelsep=period]{caption}
\graphicspath{{./}{../../docs/awes-forum-diagrams/}{../../docs/porto-2026/}}

% ── Tables ───────────────────────────────────────────────────────────────────
\usepackage{booktabs}
\usepackage{array}

% ── Math ────────────────────────────────────────────────────────────────────
\usepackage{amsmath,amssymb}

% ── Hyperlinks ───────────────────────────────────────────────────────────────
\usepackage{xcolor}
\usepackage[colorlinks=true,linkcolor=blue!60!black,citecolor=blue!40!black,urlcolor=blue!60!black]{hyperref}

% ── Custom commands ──────────────────────────────────────────────────────────
\newcommand{\RB}{$\langle$RB$\rangle$}  % Re-baseline pending marker
\newcommand{\tierP}{\textsc{[tier P]}}   % Tier P badge

% ══════════════════════════════════════════════════════════════════════════════

\title{\textbf{KiteTurbineDynamics.jl}\\[4pt]
\large Multi-Rotor Tensile Rotary Power Transmission\\for 50\,kW Airborne Wind Energy\\[8pt]
\Large Technical Report v0.2}

\author{Rod Read\\
Windswept \& Interesting Ltd, Stornoway, Scotland\\
\texttt{rod@windswept.energy}}

\date{July 2026}

\begin{document}

\maketitle
\thispagestyle{empty}

\begin{center}
\fbox{\begin{minipage}{0.92\textwidth}
\small\textbf{Status: RE-BASELINE IN PROGRESS.}
Do not cite numbers marked \RB{} until the post-settle-fix control map re-run completes (see \S0).
All results are confidence \textbf{tier P} (provisional --- single simulation model, unvalidated)
per the chart standards PRD.
\end{minipage}}
\end{center}

\vspace{12pt}
\begin{center}
\small
\textbf{Repository:} \url{https://github.com/rodread/KiteTurbineDynamics.jl} (MIT License)\\
\textbf{Tests:} full suite via \texttt{test/runtests.jl}, Julia 1.12, $\sim$3\,min full suite\\
\textbf{Intended publication:} Versioned report + Zenodo DOI (pending Gate~1)
\end{center}

\newpage
\tableofcontents
\newpage

% ══════════════════════════════════════════════════════════════════════════════
\section{Status of Numbers in This Report}
\label{sec:status}
% ══════════════════════════════════════════════════════════════════════════════

On 2026-07-04 a simulator integrity audit found that the dynamic settle stage used
$p.k_{\text{mppt}} = 614.9$ while the simulation used $k_{\text{ref}} = 15.6$ (39$\times$),
meaning \textbf{all dynamic results published before the fix started from a wrong-state settle}.
After the fix, the $\lambda = 1.0$ verification gate moved from 166 to 193\,kW (+16\%).
Consequences for this report:

\begin{itemize}
\item Every dynamic number produced before 2026-07-04 is marked \RB{} (re-baseline pending)
      and is superseded pending the control-map re-run.
\item Every result in this report is confidence \textbf{tier P} (provisional --- single simulation model,
      unvalidated) per the chart standards PRD \cite{chartprd}. There are no tier-H or tier-M results yet.
      That is not a reason to hide the badge; it is the reason to show it.
\item Two findings from the 2026-07-04 session are \textbf{new and post-fix}: the $\lambda = 0.69$
      blade-rescaled envelope pass and the $\omega$-dependent transmission-loss result (\S\ref{sec:postaudit}).
\item Earlier drafts stated 4.2$\times$ in the abstract and 4.1$\times$ in the body;
      the discrepancy is itself an argument for the consistency-stamp standard now in force.
\end{itemize}

% ══════════════════════════════════════════════════════════════════════════════
\section{Introduction: The TRPT Kite Turbine}
% ══════════════════════════════════════════════════════════════════════════════

\subsection{What is tensile rotary power transmission?}

A conventional horizontal-axis wind turbine places a rotor on a tower and transmits torque
down a rigid driveshaft. The tower must resist the overturning moment from rotor thrust---this
drives the cubic mass scaling that makes large turbines heavy.

A TRPT kite turbine \cite{tulloch2023} inverts this architecture. The rotor flies at altitude,
and torque is transmitted to the ground through a \emph{tensile} shaft: a rotating polygon of
tensioned tethers supported by thin-walled carbon-fibre rings. The rings resist the inward pull
of the tethers in hoop compression; the tethers carry tension along the shaft axis.
No tower is needed---the entire structure is held aloft by aerodynamic forces.

The shaft serves two functions simultaneously: it transmits torque from the airborne rotor to
the ground generator, and its rings provide mounting points for \emph{expansion rotors}---additional
banked rotors distributed along the shaft that contribute power and radial spreading force.

This report concerns the \textbf{power subsystem}---the TRPT shaft with expansion rotors,
operating at $\sim$30$^\circ$ elevation.
A separate \textbf{lift subsystem} (coaxial autogyro rotors on a kite line, modelled in
\texttt{CoaxialAutogyroStacking.jl}, operating at 45--55$^\circ$ elevation) provides the
axial tension that holds the TRPT aloft.
The two subsystems are modelled in separate repositories and are \textbf{not yet integrated};
the elevation mismatch is a known risk (\S\ref{sec:limitations}).

\subsection{Why multi-rotor?}

The late Prof.\ Peter Jamieson provided the foundational scaling argument for multi-rotor TRPT.
Splitting a single rotor of radius $R$ into $N$ equal stacked rotors with the same total swept
area yields mass ratio $M(1) = 0.577$ for perfectly equal rotors---a 42\% mass saving \cite{jamieson2020}.
Applied to the Windswept RodKite01 geometry, this confirmed multi-rotor TRPT as the correct
scaling direction. Note this is a geometric principle, not yet a validated TRPT result.

This insight underpins the KTD.jl architecture: rather than one large rotor,
the TRPT shaft carries multiple smaller rotors whose combined swept area produces the target
power at lower total mass.
The DE optimiser discovers the distribution of rotor sizes, ring positions, bank angles,
and tether counts simultaneously.

\subsection{The lift--power split}

The earlier 10\,kW Windswept prototype (the ``Daisy Kite'') used a static lift kite.
Dirk van Leersum's 2023 scaling analysis showed the static-kite approach becomes commercially
unviable above $\sim$10\,kW---lift kite area grows faster than power output \cite{vanleersum2023}.
The active lifting approach (coaxial autogyro stack) is intended to eliminate this bottleneck;
that claim is itself unvalidated (PCA-2 disk model, 3$\times$ solidity mismatch,
factor-of-2 error possible).

% ══════════════════════════════════════════════════════════════════════════════
\section{The KTD.jl Framework}
% ══════════════════════════════════════════════════════════════════════════════

Two computational stages:

\begin{enumerate}
\item \textbf{Differential evolution (DE) optimisation} across up to 14 design variables---60-island
      parallel population, BEM-coupled force models, multi-gate constraint system.
\item \textbf{Full 11-DoF multibody dynamic verification}---ODE solver with per-ring finite element
      analysis, inertia relief for free-floating rings, individual rotor power accounting.
\end{enumerate}

\subsection{Design variables}

\begin{itemize}
\item \textbf{Structural}: hub ring radius, ring outer diameter, wall thickness ratio ($t/D$),
      segment target length, bottom ring radius, line count
\item \textbf{Rotor}: expansion rotor count, bank angle, tip-speed-ratio scaling factors ($\lambda_i$),
      rotor radius scaling factors
\item \textbf{Constraints}: ground clearance, ring buckling FoS $\geq$ 1.5,
      collapse margin ($\delta\alpha^* - |\Delta\alpha| > 0$), bank $\leq 25^\circ$,
      tip speed $\leq$ 120\,m/s
\end{itemize}

\subsection{Verification strategy}

Every DE campaign optimum is verified by a 60-second dynamic simulation at design wind speed
(11\,m/s), computing per-ring forces, per-rotor power, shaft speed, and structural FoS at 1\,Hz.
A design that passes static gates but fails dynamic verification is flagged
\textbf{dynamically dead}---a category that includes the V10 Tight winner
(49.2\,kg static, structural failure within 60\,s of dynamic simulation).

Verification revealed the static equilibrium solver under-predicts dynamic $k_{\text{mppt}}$
by $\sim$3.3$\times$, causing DE campaigns to optimise for loads substantially lower than
dynamic operation produces.
Documented in \texttt{DECISIONS.md} \cite{ktd_decisions2026}.

\subsection{Campaign history}

\begin{table}[htbp]
\centering
\caption{DE campaign results. All targets: 50\,kW rated, 11\,m/s design wind speed.}
\label{tab:campaigns}
\begin{tabular}{@{}lrrll@{}}
\toprule
Version & Mass (kg) & $n_{\text{lines}}$ & $n_{\text{exp}}$ & Key feature \\
\midrule
V6.0 & 184.8 & 8 & 1 & Octagon baseline, 6/11 params on bounds \\
V6.2 (corrected) & 74.2 & 12 & 1 & sin formula, $\cos^{2.65}$, coupled knuckle mass \\
V9.0 & 44.5 & 8 & 9 & Dynamic $\omega$ solve, 59/60 feasible \\
V10 (unified) & 76.8 & 14 & 4 & 8 gates, unified rotor model \\
V10 Tight & 49.2$^\dagger$ & 12 & 4 & $k_{\text{mppt}} \propto \lambda^2$, ring-mapping \RB \\
V10 Reinforced & 60.8 & 12 & 4 & $+$30\% $r_{\text{bottom}}$, FoS 7.18 at 15\,m/s \RB \\
\bottomrule
\end{tabular}
\vspace{4pt}

\footnotesize
$^\dagger$Dynamically dead.
All masses \textbf{exclude expansion rotor mass} (\S\ref{sec:limitations}, item 6)
and pre-date the $m_{\text{blade}}$ $\lambda^2$ scaling fix of 2026-07-04.
The 49.2\,kg $\lambda = 1.0$ optimum was over-bladed for its ring structure.
\normalsize
\end{table}

\subsection{Known limitations}
\label{sec:limitations}

Disclosed openly because they define where collaboration is needed:

\begin{enumerate}
\item \textbf{Constant-$C_L$ expansion rotor model.} No angle-of-attack dependence,
      no induction, no Betz limit, no tip losses---the model cannot stall.
      BEM or lifting-line cross-validation is the highest-priority gap.
\item \textbf{No wake interaction between rotors.} Downstream rotors see freestream wind.
      Leuthold's axisymmetric wake induction model \cite{leuthold2019} directly addresses this.
\item \textbf{Static solver under-predicts dynamic $k_{\text{mppt}}$ by $\sim$3.3$\times$.}
      Campaigns compensate with $k_{\text{mppt\_safety}} = 3.0$; root cause (lossless torque
      propagation in the equilibrium solver) requires a dynamic-aware objective.
\item \textbf{30$^\circ$ elevation mismatch with autogyro lift} (45--55$^\circ$ sweet spot).
      Integration risk identified, not yet modelled.
\item \textbf{Blade mass is provisional} (constant-thickness skin, $\propto \lambda^2$).
      No structural blade design.
\item \textbf{Expansion rotor mass not computed}---builder sets it to zero;
      all mass totals exclude this contribution.
\item \textbf{(New, 2026-07-04)} Settle/sim $k_{\text{mppt}}$ mismatch invalidated all
      pre-fix dynamic results (\S\ref{sec:status}); five simulator integrity bugs found
      and fixed or flagged in one audit session. The audit itself is documented in
      \texttt{DECISIONS.md} and retained deliberately---tier-X history is part of the integrity story.
\end{enumerate}

% ══════════════════════════════════════════════════════════════════════════════
\section{Key Findings}
% ══════════════════════════════════════════════════════════════════════════════

\subsection{The static--dynamic gap \RB}

The most significant pre-audit finding: systematic overprediction of power by steady-state
equilibrium models versus full multibody dynamics.
AWE literature documents 18--21\% overprediction for conventional configurations
\cite{kheiri2018,leuthold2019,carceller2022}.
At the V10 Tight design point, the static solver predicted 50\,kW;
the pre-fix dynamic ODE produced 12.1\,kW---a factor of 4.1$\times$.

\textbf{\RB{} This headline number is under re-baseline.}
The 12.1\,kW dynamic figure was produced with the broken settle (\S\ref{sec:status}),
which is known to have produced partial-convergence artifacts
(the $\lambda = 1.0$ gate moved 166 $\rightarrow$ 193\,kW after the fix, $+$16\%).
The gap is expected to persist qualitatively---the two proposed mechanisms
(inter-ring torsional losses absent from the equilibrium solver, and angle-of-attack
variation under shaft twist) are real---but the magnitude (4.1$\times$ vs.\ something smaller)
must not be cited until the re-run lands.

The gap, whatever its post-fix magnitude, is not a failure of KTD.jl---it is the simulator's
primary finding. No published TRPT model has quantified it.

\subsection{Left-flank architecture}

Control-map hunting across 5--15\,m/s revealed two operating regimes for a given $k_{\text{mppt}}$:

\begin{itemize}
\item \textbf{Right flank} (high $k$, low $\omega$): high generator torque, low speed.
      Dynamic reachability challenges, poor collapse margin.
\item \textbf{Left flank} (low $k$, high $\omega$): low generator torque, high speed.
      The bisection naturally converges here.
      \RB{} The pre-fix claim ``FoS $\geq$ 2.5 while absorbing 3.4$\times$ rated capture
      at 15\,m/s'' derives from the superseded 172.7\,kW control map and awaits re-run.
\end{itemize}

Architectural decision (2026-06-30): \textbf{design for left-flank overspeed}---size blades
so that $P_{\text{min}} \leq P_{\text{rated}}$ on the left flank, size rings for thrust,
enforce tip speed $\leq$ 120\,m/s for acoustic compliance \cite{ktd_decisions2026}.

\subsection{Mass scaling}

The system-level mass scaling exponent $\alpha \approx 0.74$--$0.90$ emerges from three
competing mechanisms, not a single law:

\begin{itemize}
\item \textbf{Lines} (tension-governed): mass $\propto$ force $\propto P$. Linear.
\item \textbf{Rings} (buckling-governed, Euler column): mass $\propto P_{\text{cr}} \cdot r^2 / EI$.
      Sub-volumetric, sensitive to ring diameter.
\item \textbf{Blades} (volumetric): mass $\propto A^{3/2}$. Cubic, but blades contribute
      $<$5\% of airborne mass.
\end{itemize}

\textbf{Caveats stated plainly}: $\alpha$ is computed from \textbf{two design points}
(10\,kW, 50\,kW), excludes expansion rotor mass, and carries zero hardware validation.
The Makani precedent (designed 919\,kg, built 1,731\,kg---1.88$\times$) is the standard
we hold ourselves to: simulation-only mass claims carry limited weight until a hardware
demonstration exists.

\subsection{Post-audit findings (2026-07-04, tier P, post-settle-fix)}
\label{sec:postaudit}

\begin{enumerate}
\item \textbf{Blade-rescaled envelope pass.}
      A $\lambda = 0.69$ blade-rescaled V10 Tight
      ($k_{\text{mppt}} = 7.4$, from $k(\lambda) = k_0 \cdot \lambda^2 \cdot [r_{\text{mean}}(\lambda)/r_{\text{mean}}(1)]^3$)
      meets \textbf{P $\geq$ 50\,kW and FoS $\geq$ 1.5 across 5--15\,m/s}---168\,kW / FoS 2.51
      at 15\,m/s under unregulated MPPT (conservative), where the published V10 Tight
      failed at FoS 1.36.
      Low-wind power follows $\sim v^3$ below rated.
      This lifts the previous binding constraint.

\item \textbf{Transmission loss scales with shaft speed, not aero power.}
      Across eight runs spanning two designs and 194--303\,rpm, loss fits $c \cdot \omega^3$
      with a design-independent coefficient $c \approx 2.2$--$2.9$\,W/(rad/s)$^3$ (consistent
      with quadratic aerodynamic drag torque on the rotating shaft).
      Loss fraction rises as blades shrink (13\% $\rightarrow \sim$25\%) because generator $k$
      falls with $\lambda^2$ while shaft drag does not.
      Mechanism attribution ongoing; numerical damping (\texttt{lin\_damp}) excluded as the source.
      Low-wind deficit suggests a small additional constant-torque term; two-term fit pending.

\item \textbf{Blade-only scaling law.}
      Static aero power follows $\lambda^2$ at fixed ring radius;
      the correct $k_{\text{mppt}}$ schedule is neither $\lambda^2$ nor $\lambda^5$ alone
      but $k_0 \cdot \lambda^2 \cdot (r_{\text{mean}}\ \text{ratio})^3$, because blade radii
      are offsets from the ring radius, not the shaft axis.
\end{enumerate}

% ══════════════════════════════════════════════════════════════════════════════
\section{Comparison with Prior and Parallel Work}
% ══════════════════════════════════════════════════════════════════════════════

\subsection{University of Strathclyde---Wind Energy \& Control Centre}

The most complete published TRPT modelling framework, in close collaboration with Windswept.

\textbf{Chen et al.\ (AWEC 2026)} \cite{chen2026}: coupled aero-structural steady-state model
for TRPT power efficiency.
86.33\% transmission efficiency across 4--15\,m/s; 89\% of torque loss concentrated in the
top TRPT segment; optimal TSR shifts downward when transmission loss is considered;
stable operation requires positive tangent stiffness (minimum axial force bound).

\textbf{Amjad et al.\ (AWEC 2026)} \cite{amjad2026}: TRPT scalability framework using
QBlade LLFVW.
TRPT upscaling is transmission-limited, not rotor-limited; parametric sweeps yield a
``wide-and-short'' geometric prescription; multi-rotor identified as future work.

\textbf{Relationship to KTD.jl}: the Strathclyde models are steady-state.
Chen's 86.33\% efficiency is the best published TRPT steady-state analysis;
KTD.jl's static--dynamic gap is the natural next question---what happens when these
configurations run through full multibody dynamics?
Amjad's ``wide-and-short'' geometry is precisely what KTD.jl's DE optimiser independently produces.
Note also the convergence between Chen's transmission-loss findings and KTD.jl's $\omega^3$ loss
result (\S\ref{sec:postaudit}, item 2)---an obvious joint-validation target.

\subsection{University of Freiburg---Systems Control \& Optimization Laboratory}

\textbf{Leuthold et al.\ (AWEC 2026)} \cite{leuthold2026}: rigidly-convected lifting-line
vortex model for AWE optimal control; her 2019 wake induction model \cite{leuthold2019}
quantified 18--21\% overprediction from neglected induction.

\textbf{Diehl, De Schutter \& Harzer (AWEC 2026)} \cite{diehl2026}: vertical AWE farm
concept---99 dual-wing systems, 50\,MW on 7\,km$^2$; figure of merit is power density (MW/km$^2$).
The free-flying analogue of multi-rotor TRPT vertical stacking.
De Schutter is now at TransnetBW (grid-integration perspective).

\textbf{Relationship}: both groups distribute aerodynamic load across vertically-stacked
elements under different mechanical constraints.
KTD's fixed-geometry axisymmetric wake is a simpler validation case for Leuthold's model
than free-flying multi-kite; the static--dynamic gap is a direct data point for her research question.

\subsection{someAWE Labs---Christof Beaupoil}

The most advanced rotary AWE hardware currently flying \cite{beaupoil2026}:
autogyro pumping-mode with swashplate, active rotation compensation, Kaman-style servo flaps
for cyclic pitch (no pitch links at the hub), in cooperation with Freiburg.
Closest flying analogue to the autogyro lift subsystem; his servo-flap approach may remove
mechanical linkages between stacked units.

\subsection{Dr.\ Oliver Tulloch---TRPT inventor}

Tulloch's thesis \cite{tulloch2021} and Energies paper \cite{tulloch2023} established the
foundational TRPT model: single-section moment-to-tension ratio (MTR $\approx$ 0.05),
spring-disc dynamic formulation, $\delta\alpha^*$ collapse criterion.
All of KTD.jl's physics traces to this work.
KTD.jl extends the single-section MTR to per-section coupling from explicit ring geometry;
$\delta\alpha^*$ is implemented dynamically as \texttt{collapse\_margin} $= \delta\alpha^* - |\Delta\alpha|$,
verified healthy at 42--47$^\circ$ on the left flank.
Dr.\ Tulloch now works in offshore wind; his contribution is cited with gratitude and respect.

% ══════════════════════════════════════════════════════════════════════════════
\section{Conclusion}
% ══════════════════════════════════════════════════════════════════════════════

TRPT kite turbines with multi-rotor expansion represent a material-efficient path to
utility-relevant airborne wind energy.
KTD.jl has quantified an achievable mass range (49--61\,kg airborne at 50\,kW, tier P,
expansion-rotor mass excluded), identified the static--dynamic gap as the critical modelling
challenge \RB{}, established left-flank overspeed as the operating architecture, and---post-audit---demonstrated
a blade-rescaled design meeting power and safety criteria across the full wind envelope in simulation.

The remaining gaps---mid-fidelity aerodynamics, wake interaction, hardware validation,
dynamic-aware optimisation---map directly onto the expertise of the groups discussed in \S4.
Collaboration briefs are issued separately.
The simulator is open-source, the test suite is green, and the decisions are documented.
The next step is a conversation.

\vspace{18pt}
\begin{center}
\textit{Rod Read}\\
Windswept \& Interesting Ltd, Stornoway, Scotland\\
\texttt{rod@windswept.energy}
\end{center}

% ══════════════════════════════════════════════════════════════════════════════
% REFERENCES
% ══════════════════════════════════════════════════════════════════════════════

\begin{thebibliography}{99}

\bibitem{tulloch2023}
Tulloch, O., Yue, H., Kazemi Amiri, A.M.\ and Read, R. (2023).
A tensile rotary airborne wind energy system---modelling, analysis and improved design.
\emph{Energies}, 16(6), 2610.

\bibitem{tulloch2021}
Tulloch, O. (2021).
\emph{Tensile Rotary Power Transmission Design}.
PhD Thesis, University of Strathclyde.

\bibitem{jamieson2020}
Jamieson, P. (2020).
Top-level rotor optimisations using actuator disk theory.
\emph{Wind Energy Science}, 5, 807--818.

\bibitem{vanleersum2023}
van Leersum, D. (2023).
TRPT scaling analysis---internal report for Windswept \& Interesting Ltd.
Unpublished.

\bibitem{chen2026}
Chen, Z., Yue, H., Kazemi, A.\ and Read, R. (2026).
Power efficiency analysis of a rotary kite airborne wind energy system.
\emph{Proc.\ 11th Intl.\ Airborne Wind Energy Conference (AWEC 2026)}, Porto, Portugal.

\bibitem{amjad2026}
Amjad Zulfazli, M.M., Yue, H., Carroll, J.\ and Chen, Z. (2026).
Scalability analysis of a rotary kite AWE system with tensile rotary power transmission.
\emph{Proc.\ AWEC 2026}, Porto.

\bibitem{leuthold2019}
Leuthold, R., Crawford, C., Gros, S.\ and Diehl, M. (2019).
Engineering wake induction model for axisymmetric multi-kite AWE systems.
\emph{Wake Conference 2019}.

\bibitem{leuthold2017}
Leuthold, R., Gros, S.\ and Diehl, M. (2017).
Induction in optimal control of multiple-kite AWE systems.
\emph{IFAC-PapersOnLine}, 50(1), 153--158.

\bibitem{leuthold2026}
Leuthold, R., Crawford, C., Gros, S.\ and Diehl, M. (2026).
Trajectory tracking in AWE optimal control with a rigidly-convected, lifting-line vortex model.
\emph{Proc.\ AWEC 2026}, Porto.

\bibitem{diehl2026}
Diehl, M., Harzer, J.\ and De Schutter, J. (2026).
Vertical airborne wind energy farms based on dual-wing systems.
\emph{Proc.\ AWEC 2026}, Porto.

\bibitem{beaupoil2026}
Beaupoil, C.\ and Weyel, F. (2026).
Autogyro pumping mode rotary AWE system---design and control of a servo flap cyclic pitch actuated rotor.
\emph{Proc.\ AWEC 2026}, Porto.

\bibitem{kheiri2018}
Kheiri, M., Bourgault, F.\ and Iosilevskii, G. (2018).
Power limit and its governing factors for crosswind kite systems.
\emph{Journal of Aircraft}, 55(4), 1534--1548.

\bibitem{carceller2022}
Carceller Candau, J. (2022).
Dynamic analysis of a rotary airborne wind energy system.
MSc Thesis, Universitat Polit\`ecnica de Catalunya.

\bibitem{ktd_decisions2026}
Read, R. (2026).
\texttt{KiteTurbineDynamics.jl}---\texttt{DECISIONS.md} (architectural decision log).
\url{https://github.com/rodread/KiteTurbineDynamics.jl}

\bibitem{chartprd}
Read, R. (2026).
KTD.jl Reporting Chart Standards, PRD 0005.
\texttt{docs/prd/0005-chart-standards.md}.

\bibitem{harzer2026}
Harzer, J., Polkl\"aser, K.S., Diehl, M.\ and Moormann, D. (2026).
Multi-wing AWE research project: goals and progress.
\emph{Proc.\ AWEC 2026}, Porto.

\end{thebibliography}

\end{document}
